\section{Conclusion}
The shared execution alphabet makes menu selection and service allocation inseparable. Eligibility-aware path bounds, adaptive target covers, and guarded history aggregation give original-budget joint certificates. Harmonic pooling strengthens the reduced upper model, optimal groupwise recovery restores the original contract, and square-root cost bins improve the accuracy-dependent dimension without excluding zero costs. The method preserves every installed-level charge, accepted obligation, participation cap, and realized ceiling.

The numerical evidence separates three objects: an optimistic reduced optimum, a feasible original contract, and their certified gap. The cost-heterogeneity ablation exhibits both tighter intervals and persistent cap-driven failures; neither is hidden by a reduced solver's success. The continuous, institutional, Monge, compression, and inexact-support results retain their assumptions and proofs. The unresolved theoretical frontier is unrestricted eligibility complexity and practically small worst-case accuracy dependence, not feasibility recovery or original-budget certification. No empirical population calibration or deployment is inferred from the synthetic experiments.
